% 表紙と目次（全回版のみ）
% ============================================================
% 表紙
% ============================================================
\thispagestyle{cover}
\vspace*{18mm}
{\sffamily\color{pone}\large 名古屋市立大学 経済学部　2026年度後期　水曜3限}\par
\vspace{3mm}
{\sffamily\bfseries\fontsize{30pt}{38pt}\selectfont 基礎演習Ⅱ}\par
\vspace{2mm}
{\sffamily\fontsize{20pt}{26pt}\selectfont 各回ワークシート}\par
\vspace{4mm}
{\color{inktwo}\rule{\linewidth}{0.6pt}}\par
\vspace{4mm}
{\sffamily\small 学籍番号\makebox[40mm]{\hrulefill}\qquad 氏名\makebox[60mm]{\hrulefill}}\par
\vspace{10mm}

\begin{tcolorbox}[enhanced,colback=papertwo,colframe=papertwo,boxrule=0pt,arc=1.5mm,
    left=4mm,right=4mm,top=3mm,bottom=3mm,fontupper=\small]
{\sffamily\bfseries このワークシートの使い方}\par\vspace{1mm}
\begin{itemize}[label=・]
  \item 授業計画の各回に1枚（1〜2ページ）ずつ対応しています。授業の前に該当回を印刷するか、PDFに直接書き込んでください。
  \item 「■ ワーク」は授業内で、「次回までに」は授業外で取り組みます。
  \item AIの出力は記入欄にそのまま写さず、要点を自分の言葉で書いてください。
  \item 各回の末尾の「AI活用記録」は、提出物に添付するAI活用記録の下書きになります（書き方はAI利用ガイド）。
  \item 使用するAIは大学が提供する Copilot です。企業から提供された非公開の情報は入力しません。
\end{itemize}
\end{tcolorbox}

\vspace{6mm}
{\sffamily\bfseries 目次}\par\vspace{1mm}
{\small\makeatletter\@starttoc{toc}\makeatother}
